\documentclass{beamer}

\usetheme{Madrid}
\usepackage{booktabs}
\usepackage{graphicx}

% UJ orange - third-party approximation (brandcolorcode.com), NOT confirmed
% against UJ's own brand guidelines (uj.ac.za/brand/ blocks automated access).
% Swap this for the real value if you have the official brand PDF.
\definecolor{ujorange}{HTML}{F26522}
\setbeamercolor{structure}{fg=ujorange}
\setbeamercolor{palette primary}{bg=ujorange,fg=white}
\setbeamercolor{palette secondary}{bg=ujorange!80!black,fg=white}
\setbeamercolor{palette tertiary}{bg=ujorange!60!black,fg=white}
\setbeamercolor{title}{fg=white,bg=ujorange}
\setbeamercolor{frametitle}{fg=white,bg=ujorange}

\title[MAHER]{MAHER: Multi-Agent Hierarchical Expert\\Re-identification for Large-Scale Wildlife Recognition}
\author[Manhema et al.]{TW Manhema \and JS Phala \and HE Ashimwe \and L Phuma}
\institute[UJ]{Academy of Computer Science and Software Engineering, University of Johannesburg}
\date{\today}

% Drop the real UJ logo in as presentation/uj-logo.png (or .pdf/.jpg) and it
% shows up on every slide automatically - no other edit needed. Falls back to
% no logo if the file isn't there yet, so the deck still compiles either way.
\IfFileExists{uj-logo.png}{\logo{\includegraphics[height=0.8cm]{uj-logo.png}}}{}
\IfFileExists{uj-logo.pdf}{\logo{\includegraphics[height=0.8cm]{uj-logo.pdf}}}{}

\begin{document}

\frame{\titlepage}

% ---------------------------------------------------------------------------
\begin{frame}{Overview}
  \textbf{The problem:} identifying wild animals without physical tags is hard
  in the real world, because populations are dynamic --- new individuals keep
  appearing that a model has never seen before.

  \vspace{0.5em}
  \textbf{Our approach:} a two-stage hierarchical pipeline.
  \begin{enumerate}
    \item \textbf{Model A} (generalist) looks at a photo and predicts the
      \emph{species}.
    \item A \textbf{routing table} picks the matching \textbf{Model B}
      specialist for that species.
    \item That specialist looks at the \emph{same} photo and identifies the
      \emph{individual} animal --- or flags it as new.
  \end{enumerate}

  \vspace{0.5em}
  \textbf{Scope:} 4 species --- hyena, leopard, macaque, dog.
\end{frame}

% ---------------------------------------------------------------------------
\begin{frame}{Data \& Key Decisions}
  \begin{itemize}
    \item Source: the \texttt{wildlife-datasets} package ---
      SpottedHyenaID109, LeopardID102, MacaqueFaces, DogFaceNet.
    \item \textbf{12{,}308} labelled images across the 4 species.
    \item Scoped down from WildlifeReID-10k's $\sim$33 species to 4, for an
      honours-project timeline and CPU-only compute.
    \item \textbf{Split strategy:} individual-level, not random --- some
      individuals held out \emph{entirely} (open-set / novel), the rest keep
      most images in train but hold out a few for test too (so closed-set
      accuracy is measurable as well).
    \item Caught and fixed a real bug: an early split version left
      \emph{zero} test images belonging to a known individual, silently
      breaking closed-set evaluation.
    \item Reconciled two independently-built data pipelines (SQLite-based vs.
      a shared \texttt{metadata.csv}) that had diverged mid-project.
  \end{itemize}
\end{frame}

% ---------------------------------------------------------------------------
\begin{frame}{Model A --- Species Classifier}
  \begin{columns}
    \begin{column}{0.55\textwidth}
      \textbf{Architecture}
      \begin{itemize}
        \item ResNet50, frozen ImageNet backbone
        \item Only the final linear head trains
        \item Cheap on CPU --- fast to iterate
      \end{itemize}

      \vspace{0.5em}
      \textbf{Training}
      \begin{itemize}
        \item Cross-entropy loss, Adam optimiser
        \item 5 epochs
      \end{itemize}
    \end{column}
    \begin{column}{0.45\textwidth}
      \textbf{Result}
      \begin{block}{Test accuracy}
        \centering \Large 99.93\%
      \end{block}
      \small 4{,}069 test images --- near-perfect, since the 4 species are
      visually very distinct.
    \end{column}
  \end{columns}

  \vspace{1em}
  Two independently-trained checkpoints exist (different split
  methodologies); the deployed version uses the identity-split checkpoint
  (100\% validation accuracy).
\end{frame}

% ---------------------------------------------------------------------------
\begin{frame}{Model B --- Individual Re-identification}
  \textbf{Architecture:} ViT-Tiny backbone + ArcFace loss, trained
  \emph{separately per species} --- one specialist per species, not one
  shared flat model. This species-conditioning is the core architectural
  contribution over a flat baseline.

  \vspace{0.3em}
  \textbf{Inference:} cosine similarity against a gallery of training
  embeddings baked into each checkpoint; a similarity threshold decides
  known individual vs.\ \emph{``new animal''} (open-set rejection).

  \vspace{0.5em}
  \begin{table}
    \centering \small
    \begin{tabular}{lcccc}
      \toprule
      Species & AUC & TNR@95 & BAKS & Combined \\
      \midrule
      macaque & 0.98 & 0.94 & 0.95 & \textbf{0.94} \\
      dog     & 0.92 & 0.51 & 0.92 & 0.68 \\
      leopard & 0.74 & 0.15 & 0.65 & 0.30 \\
      hyena   & 0.72 & 0.16 & 0.51 & 0.30 \\
      \bottomrule
    \end{tabular}
  \end{table}
  \small Training data volume is the dominant factor, not species
  difficulty --- macaque/dog had far more images than hyena/leopard.
\end{frame}

% ---------------------------------------------------------------------------
\begin{frame}{UI, Integration \& Deployment}
  \textbf{Integration:} one photo in $\rightarrow$ Model A predicts species
  $\rightarrow$ routing table picks the matching Model B specialist
  $\rightarrow$ same photo $\rightarrow$ individual or \emph{``New animal''}.

  \vspace{0.5em}
  \textbf{Interfaces:}
  \begin{itemize}
    \item A native desktop app (Tkinter) for local use.
    \item A public web app (Flask) sharing the \emph{same} inference code ---
      no duplicated logic between the two.
  \end{itemize}

  \vspace{0.5em}
  \textbf{Deployment:}
  \begin{itemize}
    \item Containerised with Docker.
    \item Served behind a reverse proxy with automatic HTTPS.
    \item Live, publicly accessible demo.
  \end{itemize}

  \vspace{1em}
  \centering \large \textbf{Live demo:} \texttt{\small https://maher.lwazi.gay}
\end{frame}

\end{document}
